\begin{abstract}
When tests fail, developers increasingly hand the failure to an LLM and ask where the bug is. Research tools propose
far more elaborate pipelines that claim to ground the model's reasoning in program execution. We test that claim on
SemLoc, a recent localizer for semantic faults that turns LLM-inferred constraints into a violation spectrum, and
compare it with the everyday practice. We recompute every published configuration from the authors' logs, audit the
evaluation, check the benchmark's labels by execution and by review, rerun everything three times with a current model,
and ablate each signal. Four findings stand out. First, SemLoc's top-1 answer is guessed by design: an undocumented bonus
always ranks the LLM's own guess first, and the counterfactual step never changes a ranking. Second,
\LabaudAllMoved{} of the benchmark's \LabaudAllN{} fault labels move under review (\LabaudAllBlank{} sit on blank lines),
and the mislabels happen to favor SemLoc. Third, a single prompt with the failing tests beats the pipeline at Acc@1 on all
three label sets (\RqTwoRepairedAllLlmTestsDAccOne{} points on execution-repaired labels) at about a seventh of the
tokens; ablation shows the failing tests, not per-line coverage, carry the signal. Fourth, nothing more elaborate
helps---adding coverage to SemLoc, verifying candidate lines with test-passing patches, voting, or static analysis with
ruff and pyrefly---because a small patch passes all tests at \PavmechReviewedOtherLineFixable\% of non-faulty lines:
the tests suggest where the bug is but cannot prove it. Code, data, LLM logs and a replayable cache are public.
\end{abstract}
